% ! TeX root = ../main.tex

\section{Background and Motivation}\label{sec:orca:bg}

\emph{Bulk synchronous parallel} (BSP)~\cite{Valia1990:BSP} workloads in both scientific computing and machine learning
training pose unique efficiency challenges.
These jobs, increasingly spanning tens of thousands of nodes, execute in lockstep and synchronize
frequently via collective communication.
At these scales, BSP is uniquely fragile: a single straggler can block cluster-wide progress at
synchronization points.
Their root causes are highly context-dependent, ranging from load imbalance and suboptimal overlap
to tuning-induced variability, and have been documented
extensively~\cite{\citeGroupVarHpc,\citeGroupVarAmr,\citeGroupVarRdma,\citeGroupVarMisc,\citeGroupML}.
Despite extensive case studies, performance pathologies persist and remain costly to diagnose and
fix.

We argue this reflects a structural limitation in how BSP observability is practiced.
While monitoring infrastructure is widely deployed~\cite{:Prometheus,Schwa2021:LDMS}, it targets
\emph{known unknowns} via a fixed set of metrics such as CPU utilization.
Performance pathologies are dominated by \emph{unknown
  unknowns}~\cite{Major2017:o11y,:OtelO11y,Sridh:DistributedSystemsObservability} ---
they cannot be anticipated by predefined metrics, and resolving them requires iterative
investigation at varying resolutions: from high-level aggregates to fine-grained detail when
needed.
Current approaches to capturing fine-grained telemetry --- large event traces followed by
post-mortem analyses --- incur prohibitive collection and analysis
costs~\cite{Yildi2023:IOMax,\citeJainAmr}, delaying interventions.
As a result, operators are forced to choose between expensive investigations and ad-hoc
workarounds.
This gap motivates a different model: always-on application observability as a production
capability, with relevant views constructed on demand in real-time, hypotheses tested live, and
collection steered accordingly while the job is still running --- a paradigm we call
\emph{steerable observability}.

We now illustrate this gap with case studies from HPC and ML training (\S\ref{sec:orca:bg:studies}),
survey existing approaches and their limitations (\S\ref{sec:orca:bg:gaps}), and use our experiences
with eBPF~\cite{:eBPF} to motivate a programmable substrate for real-time visibility at workload
scale (\S\ref{sec:orca:bg:missing}).

% This gap motivates a different model: always-on observability as a production capability, with
% relevant views derived in real-time, hypotheses tested live, and collection steered accordingly
% while the job is still running --- a paradigm we call \emph{steerable observability}.
% Performance pathologies are dominated by \emph{unknown
%   unknowns}~\cite{Major2017:o11y,:OtelO11y,Sridh:DistributedSystemsObservability} ---
% resolving them requires constructing the right views from telemetry and refining them until causal
% attribution is possible.
% Dominant paradigms make this loop slow and rare: diagnosing pathologies requires rerunning
% applications under progressively heavier diagnostics, from lightweight profiles to fine-grained
% per-rank traces, followed by expensive \emph{extract-transform-load} (ETL)~\cite{Yildi2023:IOMax}
% before queries can be served.
% As a result, operators are forced to choose between expensive investigations and ad-hoc
% workarounds.
% This gap motivates a different model: always-on observability as an operational capability, with
% relevant views derived in real-time, hypotheses tested live, and collection steered accordingly
% while the job is still running.


\subsection{Case Studies from HPC and ML}\label{sec:orca:bg:studies}

\subsubsection{Optimizing Placement in AMR}
\figBgVolatile{}

% Our study of placement and load-balancing in \emph{adaptive mesh refinement} (AMR)
% codes~\cite{\citeJainAmr} illustrates this gap concretely.
A prior study of placement and load-balancing in \emph{adaptive mesh refinement} (AMR)
codes~\cite{\citeJainAmr} illustrates this gap concretely.
While widely known for load-imbalance challenges~\cite{Dubey2014:Survey}, the study finds
load-imbalance indistinguishable from other hardware, software, and tuning artifacts without
fine-grained telemetry.
\Cref{fig:orca:bg-volatile} shows one such example: tail latency spikes in \mpiw{}
that only manifested at scale would get attributed to collectives in aggregate telemetry.
Diagnosing each artifact required months of iterative investigation, while mitigations typically took days.

The primary diagnostic bottleneck was a slow feedback loop: each iteration required running the
codes, collecting large event traces, followed by post-mortem trace analyses bottlenecked on
expensive ETL (\emph{extract-transform-load}) from formats such as CSV.
Mitigations identified were cluster- and workload-specific.
Even after these artifacts were fixed, placement demonstrated a fundamental tradeoff
between compute load balance and communication locality --- the right balance between the two
varies across workloads, scales, and architectures.
Translating these optimizations to production environments would require repeating this
time-consuming exercise --- underscoring the need for infrastructure that makes such investigations
routine.
% Translating them to production environments would require repeating the same exercise.
% Given that optimal configurations are not portable across environments, the need is not for better
% one-off analyses but for infrastructure that makes such investigations routine.

\subsubsection{Parallels from the Pytorch Ecosystem}

The PyTorch~\cite{Paszk2019:PyTorch} ML training ecosystem exhibits the same offline feedback loop.
Kineto~\cite{2025:PytorchKineto} captures fine-grained traces on demand as JSON, triggered via a
node-local control daemon called Dynolog~\cite{2025:MetaDynolog}.
\emph{Holistic Trace Analysis} (HTA)~\cite{2025:MetaHTA} computes derived views such as activity
breakdowns and straggler detection.
Together, these cover capture, actuation, and analysis --- but HTA is a single-node
Pandas~\cite{:Pandas} script that cannot scale to cluster-wide traces, and Dynolog provides
best-effort actuation with no workload-level consistency guarantees.
% , whereas reliable construction of views like
% \Cref{fig:orca:bg-volatile} requires complete telemetry for each collective instance.

Recent ML training systems have begun addressing specific pathologies online --- straggler
localization~\cite{Yao2025:Holmes,Wu2025:GREYHOUND,\citeCuiXputimer,Guan2025:PerfTracker,Deng2025:Mycroft},
fault diagnosis~\cite{Dong2025:Aegis,Deng2025:Minder}, and training
correctness~\cite{Jiang2025:TrainCheck} --- but each reimplements subsets of collection,
aggregation, and control from scratch.
A common programmable substrate would generalize these diagnostics, allowing detection strategies
to be expressed as analytical queries over system state, composable without
rebuilding the underlying infrastructure.

\subsection{Existing Approaches and Gaps}\label{sec:orca:bg:gaps}

\parahead{Monitoring}
Monitoring enables basic cluster visibility via periodically captured metrics, but lacks the
workload context necessary to construct views like \Cref{fig:orca:bg-volatile}.
Widely deployed infrastructure for metrics includes Prometheus~\cite{:Prometheus}, a time-series
database, and LDMS~\cite{Schwa2021:LDMS}, an on-fabric hierarchical metric collection
infrastructure used in HPC clusters.

\parahead{Profiling and Tracing}
Tools such as TAU~\cite{Shend2006:TAU}, \scorep{}~\cite{Knupf2012:Score-P}, and
Kineto~\cite{2025:PytorchKineto} can provide both high-level aggregates via profiles or fine-grained telemetry
via traces, but are post hoc by design.
Traces can characterize short runs, but cannot provide always-on visibility over long-running jobs.
In the absence of any real-time feedback to guide collection, traces simultaneously overcollect
and undercollect --- despite capturing every \mpiw{}, they can not explain which sub-functions
within anomalous \mpiw{} caused the latency spikes.
Formats like OTF2~\cite{Eschw2012:OTF2} and JSON~\cite{:ChromeJsonTEF} 
require expensive transforms for dataframe-style analytics~\cite{Jain2025:AMR,Yildi2023:IOMax},
and visual inspection for anomalies is impractical for cluster-scale workloads.

% \cite{Bhate2019:Hatchet} (Hatchet)
% \cite{Yildi2023:IOMax} -- DFTracer cousin, exists because they do not know Polars + Lazy Parquet exists
\parahead{Dataframe Analytics}
Recent tools such as Caliper~\cite{Boehm2016:Caliper}, \dftracer{}~\cite{Devar2024:DFTracer},
Pipit~\cite{Bhate2024:Pipit}, and Pytorch HTA~\cite{2025:MetaHTA} have adopted
Pandas~\cite{:Pandas} and Dask-based~\cite{Rockl2015:Dask} dataframe analytics interfaces to
analyze telemetry, but require expensive ETL~\cite{Yildi2023:IOMax,\citeJainAmr} 
--- the working set for which scales with \emph{ranks $\times$ timesteps}.
%formats such as CSV and JSON 
Columnar formats like Parquet~\cite{:ApacheParquet} offer efficient post-hoc analytics, but do not
by themselves address overcollection, working set scaling, or real-time feedback, and remain
largely unadopted in tracing workflows.
In a simple internal port of HTA analyses from JSON to Parquet, we observed 12$\times$ faster data
loading and 8$\times$ faster end-to-end analysis.

% \cite{Godoy2020:ADIOS2} (ADIOS)
% \cite{Eisen2025:ADIOS} -- ADIOS SST
\parahead{Streaming and In-situ Analytics}
ADIOS~\cite{Godoy2020:ADIOS2} enables fixed-function in-situ workflows~\cite{Eisen2025:ADIOS} on
scientific data, typically run as another MPI application.
%provides transport abstractions for in-situ
MRNet~\cite{Roth2003:MRNet,Willi2016:MRNetApps} is a tree-based overlay used as a reduction backend
for commercial debuggers, but provides a compiled packet-based programming model.
Neither enable runtime-programmable dataframe-style analytics needed for real-time BSP
observability.
Cloud streaming systems such as Flink~\cite{Carbo2015:ApacheFlink} and Kafka~\cite{Kreps2011:Kafka}
are designed for independent event streams over IP networks They lack fabric-aware transport and
BSP-aligned coordinated collection, precluding views like \Cref{fig:orca:bg-volatile}.

\subsection{The Missing Substrate}\label{sec:orca:bg:missing}

Diagnosing the root cause behind \Cref{fig:orca:bg-volatile} required capabilities beyond what any
single tool in the taxonomy provides.
At node scale, eBPF~\cite{:eBPF} provides a template: dynamic probes, programmatic aggregation, low
overhead in the common case with depth when needed.
It enabled controlled and targeted access to telemetry within \mpiw{} after alternatives were
exhausted.
Hypotheses about interrupts and scheduler effects were tested and discarded in hours.
But eBPF is node-local, and our workflow involved ad-hoc orchestration, uncoordinated collection,
and manual post-hoc correlation.
Extending this to workload scale requires a substrate for real-time \emph{feedback}, treating
observability not as distributed logging, but as materializing views directly from running
applications.
But visibility alone is insufficient for steerability. Interventions --- whether additional views
or parameter adjustments --- require a \emph{control plane} back to the running application,
and may be driven by humans, automated agents, or policies.
The next section states the requirements such a substrate must satisfy.
